\section{Tube-Distance Concentration}
\label{app:tube-tails}
This result isolates the probability of leaving a declared tube from the
geometric distance estimate used to define that tube. In particular, a
positive reach condition alone is not a derivation of a residual bound.

\begin{proposition}[Tube-distance concentration from residual tails]
\label{prop:tube-concentration}
Let $\Gamma_b=\gamma_b(\mathcal S_b)$ and
$d_{\Gamma_b}(x)=\inf_{y\in\Gamma_b}\|x-y\|$. On a measurable coordinate
subset $U_b$ of positive branch-coordinate probability, suppose the observation is
$X=\gamma_b(S)+t_b(S)\varepsilon_{\parallel,b}
+N_b(S)\varepsilon_{\perp,b}$. Let
$0<M_{N,b}=\sup_{s\in U_b}\|N_b(s)\|_{\mathrm{op}}<\infty$.
Assume a nonnegative measurable remainder $R_b$ and radius $\rho_b>0$
satisfy
\[
 d_{\Gamma_b}(X)\le M_{N,b}\|\varepsilon_{\perp,b}\|+R_b
 \quad\text{whenever }S\in U_b\text{ and }
 M_{N,b}\|\varepsilon_{\perp,b}\|+R_b<\rho_b.
\]
For $0<r<\rho_b$ and $0<\eta<1$, set $a=(1-\eta)r/M_{N,b}$ and
\begin{align*}
 \Delta_b(r,\eta)
 &=\Pbb(S\notin U_b\mid Z=b)\\
 &\quad+\operatorname*{ess\,sup}_{s\in U_b}
          \Pbb(\|\varepsilon_{\perp,b}\|>a\mid S=s,Z=b)\\
 &\quad+\operatorname*{ess\,sup}_{s\in U_b}
          \Pbb(R_b>\eta r\mid S=s,Z=b).
\end{align*}
Essential suprema use the branch-coordinate law. Then
\[
 \Pbb\{d_{\Gamma_b}(X)>r\mid Z=b\}\le\Delta_b(r,\eta).
\]
For $n_b$ observations conditional on branch $b$, the probability that at
least one leaves the radius-$r$ tube is at most $n_b\Delta_b(r,\eta)$.
For $n$ draws from the mixture and an active branch subset $\mathcal B_0$,
\[
 \Pbb\{\exists i:Z_i\in\mathcal B_0,
                  d_{\Gamma_{Z_i}}(X_i)>r\}
       \le n\sum_{b\in\mathcal B_0}\pi_b\Delta_b(r,\eta).
\]
\end{proposition}
\begin{proof}
If $S\in U_b$, $\|\varepsilon_{\perp,b}\|\le a$ and $R_b\le\eta r$,
the assumed bound is at most $r<\rho_b$. The escape event is therefore
contained in the union of the three complement events in $\Delta_b$.
Conditioning on $S$ and applying the union bound proves the first assertion.
Union bounds over observations and then mixture labels give the other two.
They do not require independence beyond the stated one-observation laws.
\end{proof}

\begin{corollary}[Gaussian residual-tail specialization]
\label{cor:gaussian-tube-tail}
In Proposition~\ref{prop:tube-concentration}, suppose conditionally on $S=s$
that $\varepsilon_{\perp,b}=A_b(s)G$ with
$G\sim\Normal(0,I_{r_b^{\varepsilon}})$. Set
\[
 \bar m_b=\operatorname*{ess\,sup}_{s\in U_b}\E\|A_b(s)G\|,
 \qquad
 \bar\sigma_b=\operatorname*{ess\,sup}_{s\in U_b}\|A_b(s)\|_{\mathrm{op}}.
\]
If $0<\bar\sigma_b<\infty$, the normal-residual term in $\Delta_b$ is
bounded by
\[
 \exp\left\{-\frac{(a-\bar m_b)_+^2}{2\bar\sigma_b^2}\right\}.
\]
If $\bar\sigma_b=0$, that term is zero for $a>0$.
\end{corollary}
\begin{proof}
For fixed $s$, the Gaussian function $g\mapsto\|A_b(s)g\|$ is
$\|A_b(s)\|_{\mathrm{op}}$-Lipschitz. Gaussian Lipschitz concentration
\citep{boucheron2013concentration} bounds its upper tail above its mean
by $\exp\{-t^2/(2\|A_b(s)\|_{\mathrm{op}}^2)\}$ for $t\ge0$.
Use the trivial bound one when $a$ is below the mean, otherwise take
$t=a-\E\|A_b(s)G\|$. Uniform bounds on the mean and Lipschitz constant
give the display. Zero operator norm gives zero residual almost surely.
\end{proof}

Other observation kernels enter the same first proposition through their
own residual tails. It is the kernel law and the assumed geometric remainder,
not the name of the mixture family, that determine concentration.
